\section{Further Robustness}
\label{app:online}

This appendix states sufficient conditions for the two-country sign result of \cref{rem:robustness} and relaxes two maintained assumptions of the nested-CES model, linear technology and balanced trade. Throughout, $e_{Aj}$ is units of $A$'s currency per unit of $j$'s currency, so a rise is a depreciation of $A$ against $j$.

\subsection{Two-country sufficient conditions}
\label{app:online-twocountry}

Let $e\equiv e_{AB}$ and let $\tau$ be $A$'s tariff on imports from $B$. Border prices are one in the producer's currency, so with import and export quantities $Q_M$ and $Q_X$, $A$'s import payments are $M_A=eQ_M$, its export receipts are $X_A=Q_X$, and
\begin{equation*}
\TB_A(e,\tau)\equiv X_A(e,\tau)-M_A(e,\tau),
\end{equation*}
where $X_A$ and $M_A$ include all induced effects, including income effects of tariff revenue. If $\TB_A$ is continuously differentiable with $\TB_{A,e}\neq0$ at an equilibrium, the implicit function theorem gives a locally unique $e(\tau)$ with
\begin{equation*}
\frac{de}{d\tau}=-\frac{\TB_{A,\tau}}{\TB_{A,e}}.
\end{equation*}

The first condition is $\TB_{A,\tau}>0$: at fixed $e$, the tariff-induced fall in import payments exceeds any tariff-induced fall in export receipts. Because $M_A=eQ_M$ is measured at border prices, import payments fall if and only if the import quantity falls; the condition does not require import expenditure $(1+\tau)eQ_M$ to fall. Under a lump-sum rebate of revenue, the income effect can offset substitution. With $\varepsilon^H$ the compensated own-price elasticity of import demand, $\eta_I$ its income elasticity and $\theta$ the revenue share of income,
\begin{equation*}
\frac{d\log Q_M}{d\log(1+\tau)}\bigg|_e=-\frac{\varepsilon^H}{1-\eta_I\theta},
\end{equation*}
so at free trade ($\theta=0$) the import quantity falls if and only if $\varepsilon^H>0$. In the nested-CES model at free trade, $\varepsilon^H=\rho_D>0$, the domestic diversion margin of \cref{eq:twocountry}, and $X_{A,\tau}=0$ because $B$'s income and demand are unchanged at a fixed exchange rate.

The second condition is $\TB_{A,e}>0$. Since $eX_{A,e}=X_A\varepsilon_X$ and $eM_{A,e}=M_A(1-\varepsilon_M)$, with $\varepsilon_X$ and $\varepsilon_M$ the quantity elasticities of \cref{sec:bilateral}, at a balanced baseline $e\,\TB_{A,e}=X_A(\varepsilon_X+\varepsilon_M-1)$, and the condition is equivalent to the Marshall--Lerner condition $\varepsilon_X+\varepsilon_M>1$.

When both conditions hold, $de/d\tau<0$: the tariff appreciates $A$'s currency against $B$. These are sufficient conditions, not properties of every model: the fiscal closure (revenue wasted, spent by the government or transferred abroad) and the pricing convention (local- or dominant-currency pricing, incomplete pass-through, markups) change the quantity responses and can reverse either sign.

\subsection{Sector-specific factors}
\label{app:online-rv}

Replace the linear tradable technology with a Ricardo--Viner technology, $Y_{Ti}=A_{Ti}L_{Ti}^{\gamma}$ with $0<\gamma\leq1$, where a fixed sector-specific factor owned by domestic households earns rents $(1-\gamma)Y_{Ti}$; non-tradables remain linear in labor, which is mobile across sectors. At a free-trade baseline, tradable employment is $L_{Ti}=\gamma\alpha_TL_i/(1-\alpha_T+\gamma\alpha_T)$ whatever the exchange rates, so the adjustment matrix is that of the linear model with the same nominal incomes. Only the disturbance changes, through the tariff-setter's income: the rebate draws labor out of tradables and raises the wage, so with $s_{AB}$ the share of $A$'s income spent on $B$'s good,
\begin{equation*}
\frac{d\log I_A}{d\tau_{AB}}\bigg|_{\nu}=\kappa(\gamma)\,s_{AB},\qquad
\kappa(\gamma)\equiv1+\frac{(1-\alpha_T)^2(1-\gamma)}{\alpha_T\,(1-\alpha_T+\gamma\alpha_T)}\geq1,
\end{equation*}
with equality if and only if $\gamma=1$. \Cref{prop:unilateral} then holds with $\rho_D$ replaced by
\begin{equation*}
\rho_D(\gamma)\equiv\rho_D-(\kappa(\gamma)-1)m=\alpha_D\eta+(1-\alpha_D)(1-\alpha_T)\,\frac{\gamma}{1-\alpha_T+\gamma\alpha_T}\leq\rho_D,
\end{equation*}
so the bystander threshold falls to $N\rho_D(\gamma)$, while the NEER still appreciates because $\rho_D(\gamma)>0$. At the US-like parameters and $N=3$, the threshold is $3.96$ with linear technology and $3.83$ at $\gamma=1/2$.

\subsection{Permanent trade imbalances}
\label{app:online-imbalance}

Let $T_i$ be permanent transfers fixed in $A$'s currency, positive for a receipt, with $\sum_{i\in\mathcal I}T_i=0$. Each country's external-balance condition becomes
\begin{equation*}
\TB_i+T_i=0\qquad\text{for every }i\in\mathcal I,
\end{equation*}
so a recipient runs a trade deficit equal to its receipt. The general linear system \cref{eq:response} remains valid at any such baseline.

As an illustration of U.S. scale, let $A$ receive a transfer equal to 3 percent of its income, paid by the target $B$, with the US-like parameters and $N=3$. Because $A$'s import share is small, financing the deficit requires a baseline appreciation of $A$ of about $15$ log points against $B$. The nominal bystander threshold rises from $3.96$ to $4.94$: $A$'s imports tilt toward the cheapened target, so the bystander is less exposed to diversion from $A$'s tariff. The single-threshold structure survives on $\rho\in[1.05,60]$. Alternative transfer allocations and tariff levels are reported in the replication package.
